\section{Complete a small rectangle before moving on}
A tiled transpose divides the source into small rectangles and finishes one
rectangle at a time. This exercise uses a tile side of four. The tile whose
upper-left coordinate is \((i_0,j_0)\) contains
\[
  i_0\leq i<\min(i_0+4,R),\qquad
  j_0\leq j<\min(j_0+4,C).
\]
Its first column is \(j_0\), not zero. Its first row is \(i_0\), not zero.
The tile origin advances in steps of four; the cells inside a tile advance in
steps of one.

The minimum operations clip a tile at the matrix edges. In a three-by-five
matrix, the tile at \((0,0)\) contains three rows and four columns. The tile
at \((0,4)\) contains three rows and one column. Both are smaller than a full
four-by-four tile. Neither may read padding as if it were part of the matrix.

\begin{MARSWorkedExample}{Two tiles cover fifteen cells}
  The first tile visits source indices
  \(0,1,2,3,5,6,7,8,10,11,12,13\) when its internal traversal is row-major.
  The second visits \(4,9,14\). Compare this with the ordinary source order,
  which visits \(0,1,2,3,4,5,\ldots,14\). Both cover the same cells, and both
  apply the same destination formula. Their order differs before the fifth load.
  That point makes the small rectangular example useful for checking that tiling
  has actually been implemented.
\end{MARSWorkedExample}

For a tile row starting at source row \(i\), initialize its source pointer to
\(\operatorname{input}+4(iC+j_0)\), and its destination pointer to
\(\operatorname{output}+4(j_0R+i)\). The inner cell loop then uses the same
four-byte source advance and \(4R\)-byte destination advance as the ordinary
version. Recompute the starting pointers at the next tile row. This is a useful
first implementation because the formulas remain visible at each row boundary.
Later, pointer recurrences can reduce those multiplications if measurement
justifies the added state.

Every valid source coordinate belongs to exactly one tile. Integer division of
its row and column by four gives the tile's row and column numbers. Those
quotients are unique, and the clipped bounds retain every coordinate that lies
inside the matrix. Finishing every tile therefore retains the ordinary
version's coverage and one-to-one destination mapping.

Tiles may be visited in different orders. One implementation finishes all tile
columns in the first tile row before starting the next tile row. Another finishes
all tile rows in the first tile column before moving right. Both satisfy this
exercise's tiled contract. Register choices and address recomputation strategies
can differ as well. The observable requirement is to complete all reads and
writes for one clipped tile before accessing another tile, with one aligned
word read and one aligned word write for each matrix cell. A read must precede
its corresponding write. This permits buffering several words in registers
within the tile.

The ordinary contract instead requires the source reads and corresponding output
writes to each follow source row-major order. It does not require the load and
store for a cell to be adjacent instructions. These traversal contracts make
the comparison an exercise in two memory access patterns. Equal final output
alone would accept two copies of the ordinary algorithm and would not establish
the tiled objective.
